\documentclass[11pt,a4paper]{article}
\usepackage[T1]{fontenc}
\usepackage{lmodern}
\usepackage{amsmath}
\usepackage{amssymb}
\usepackage[table]{xcolor}
\usepackage{tabularx}
\usepackage{multirow}
\usepackage{geometry}
\usepackage{enumitem}
\usepackage{parskip}
\usepackage{titlesec}
\usepackage{multicol}
\usepackage{fancyhdr}
\usepackage{needspace}
\usepackage[most]{tcolorbox}
\usepackage[colorlinks=true,linkcolor=blue!50!black,urlcolor=blue!50!black,linktoc=none]{hyperref}

\geometry{margin=2.5cm}

% Short forced heading lines (sub-answer labels, topic headers) are intentionally
% short; stop TeX reporting them as underfull (no overfull boxes exist).
\hbadness=10000
\vbadness=10000

% ===== Pastel colours (cycled per question number) =====
\definecolor{q1color}{HTML}{D6EAF8}
\definecolor{q2color}{HTML}{D5F5E3}

\titleformat{\section}{\large\bfseries}{}{0em}{}

% Question heading: unnumbered, never orphaned at a page bottom
\newcommand{\qsection}[1]{\needspace{10\baselineskip}\section*{#1}}

% Running headers: current session left, page number right
\pagestyle{fancy}\fancyhf{}
\fancyhead[L]{\small\itshape\leftmark}
\fancyhead[R]{\small\thepage}
\renewcommand{\headrulewidth}{0.3pt}

% Mark-scheme table (Paper 2 only), tinted per question colour
\newenvironment{mstab}[1]{%
  {\color{#1!75!black} Mark Scheme}\par\nopagebreak\smallskip
  \arrayrulecolor{#1!70!black}%
  \renewcommand{\arraystretch}{1.3}%
  \tabularx{\textwidth}{%
 |>{\columncolor{#1!12}\raggedright\arraybackslash}p{1.9cm}|%
 >{\columncolor{#1!12}}X|%
 >{\columncolor{#1!12}\centering\arraybackslash}p{1.1cm}|}%
  \hline
  \rowcolor{#1!45} Question & Answer & Marks \\ \hline
}{\endtabularx\arrayrulecolor{black}\par\medskip}

% Exemplar-answer box
\newtcolorbox{ansbox}[1]{
  colback=#1!8, colbacktitle=#1!30, colframe=#1!70!black, coltitle=black,
  fonttitle=\normalfont, title=Exemplar Key, boxrule=0.6pt, arc=2mm,
  left=4pt, right=4pt, top=4pt, bottom=4pt, breakable}

% Session banner: flows inline, no page break, sets the running header
\newcommand{\sessionheader}[3]{%
  \par\Needspace*{14\baselineskip}\bigskip
  \markboth{#1 - #2}{}%
  \begin{tcolorbox}[colback=black!6, colframe=black!55, arc=2mm,
 boxrule=0.9pt, left=8pt, right=8pt, top=7pt, bottom=7pt]
 \centering {\LARGE\bfseries #1 \ - \ #2}\\[5pt] {\normalsize #3}
  \end{tcolorbox}\medskip}

\begin{document}

\sessionheader{5054/22/M/J/25}{May/June 2025}{Theory: Mark Scheme \& Model Answers}

\label{sess:p2_mj25}

\qsection{5054/22/M/J/25 - Q1}\label{p2_mj25:q1}

\begin{mstab}{q1color}
 & points plotted to within $1/2$ a small square & B1 \\ \cline{2-3}
\multirow{-2}{1.9cm}{1(a)(i)} & smooth curve drawn & B1 \\ \hline
 & (average speed =) $d/t$ in any form & C1 \\ \cline{2-3}
\multirow{-2}{1.9cm}{1(a)(ii)} & 195 (cm/s) & A1 \\ \hline
 & draw a tangent (at $t = 0.40$ s) & B1 \\ \cline{2-3}
\multirow{-2}{1.9cm}{1(a)(iii)} & find the gradient / slope (of the graph) \newline or the gradient of the tangent & B1 \\ \hline
 & mention of air resistance / resistive force & B1 \\ \cline{2-3}
 & resultant force decreases \newline or air resistance / upwards force increases (as speed increases) / becomes high(er) & B1 \\ \cline{2-3}
\multirow{-3}{1.9cm}{1(b)(i)} & (eventually) air resistance = weight (and acceleration is zero) \newline or air resistance balances / cancels weight & B1 \\ \hline
 & straight line (of positive) \newline or (line of) constant gradient / slope & B1 \\ \cline{2-3}
\multirow{-2}{1.9cm}{1(b)(ii)} &  &  \\ \hline
\end{mstab}

\begin{ansbox}{q1color}
(a)(i) Plotting the Graph\\[4pt]
Plot each coordinate pair from Table 1.1 carefully on the grid, ensuring points are within half a small square of their exact positions. Then, draw a single smooth curve passing through these plotted points.

\vspace{4pt}
(a)(ii) Average Speed\\[4pt]
\[ \text{average speed} = \frac{\text{total distance}}{\text{total time}} \]
At $t = 0.40 \text{ s}$, the distance fallen is $78 \text{ cm}$.
\[ v = \frac{78}{0.40} \]
\[ \text{speed} = \boxed{195 \text{ cm/s}} \]

\vspace{4pt}
(a)(iii) Instantaneous Speed\\[4pt]
To find the speed at exactly $t = 0.40 \text{ s}$, draw a straight tangent line touching the curve at the point where $t = 0.40 \text{ s}$. Then, calculate the gradient (slope) of this tangent line.
\clearpage
\vspace{4pt}
(b)(i) Changing Acceleration\\[4pt]
As the ball falls and its speed increases, the upward force of air resistance acting against it also increases. This increasing upward force opposes the constant downward weight of the ball, meaning the overall downward resultant force decreases (causing acceleration to decrease). Eventually, the upward air resistance fully balances the downward weight, making the resultant force zero, which means the acceleration becomes zero.

\vspace{4pt}
(b)(ii) Graph with Zero Acceleration\\[4pt]
When acceleration is zero, the ball moves at a constant terminal velocity, so the distance-time graph becomes a straight line with a constant positive gradient.
\end{ansbox}

\end{document}
